\section{Origins: the wiring experiments that motivated front wheeler}
\label{sec:exp1}

\fw{} rests on a premise: most per-step improvement does not require deep
credit assignment. The experiments in this section---originally aimed at a
different question---supplied that premise. Supermask results show that a
randomly weighted, sufficiently wide network already contains subnetworks
that solve the task \citep{ramanujan2020whats, zhou2019deconstructing,
malach2020proving}; we asked how \emph{local} the search for such a
subnetwork can be made, and the control conditions of that study are what
pointed at \fw{}.

\paragraph{Setup.} MLPs with frozen random weights (signed-constant
initialization); only a binary mask (top-$k$ scores per layer, $50\%$ kept)
is learned. Rules, ordered by the information available to hidden layers:
\ep{} (straight-through gradients on scores, i.e.\ full BP---the upper
bound), \dfa{} (the layer's error is a fixed random projection of the output
error \citep{nokland2016direct}, applied here to mask scores: no weight
transport, no chained backward), and \rmhebb{} (hidden layers see one global
scalar per sample modulating a batch-centered Hebbian eligibility
\citep{fremaux2016neuromodulated, williams1992simple}). Controls:
\textsc{random} mask, \textsc{readout} (train the final layer only), and
\textsc{dense} (ordinary training). 30/40 epochs (MNIST/CIFAR-10), 3 seeds.

\begin{table}[t]
\centering
\caption{Final test accuracy (mean $\pm$ std, 3 seeds), at two and four
hidden layers. Only the mask is learned in \ep{}/\dfa{}/\rmhebb{}.}
\label{tab:exp1}
\begin{tabular}{lcccc}
\toprule
& \multicolumn{2}{c}{2 hidden layers} & \multicolumn{2}{c}{4 hidden layers} \\
Rule & MNIST & CIFAR-10 & MNIST & CIFAR-10 \\
\midrule
\textsc{dense} (control) & $98.7 \pm 0.1$ & $53.0 \pm 0.4$ & --- & --- \\
\ep{} & $98.2 \pm 0.1$ & $\mathbf{55.9 \pm 0.2}$ & $98.3 \pm 0.1$ & $\mathbf{56.2 \pm 0.1}$ \\
\dfa{} & $96.9 \pm 0.1$ & $53.3 \pm 0.2$ & $81.8 \pm 0.9$ & $23.1 \pm 8.9$ \\
\textsc{readout} (control) & $94.9 \pm 0.1$ & $45.6 \pm 0.2$ & --- & --- \\
\rmhebb{} & $51.7 \pm 3.9$ & $26.3 \pm 6.3$ & $34.1 \pm 9.9$ & $18.2 \pm 4.0$ \\
\textsc{random} (control) & $12.5 \pm 3.5$ & $10.1 \pm 0.8$ & --- & --- \\
\bottomrule
\end{tabular}
\end{table}

\paragraph{What holds, and what is shallow-only.} Three results
(Table~\ref{tab:exp1}, Figure~\ref{fig:exp1curves}). First, \emph{wiring
rivals weights, robustly}: gradient mask search over frozen random weights
comes within half a point of dense training on MNIST, beats it on CIFAR-10
($55.9\%$ vs.\ $53.0\%$, where the dense MLP overfits), and is
depth-stable---at four hidden layers it is unchanged or better. Second,
\emph{the local rule's parity is a shallow-network phenomenon}: at two
hidden layers \dfa{} mask search matches dense weight training on CIFAR-10,
but at four hidden layers it collapses ($96.9 \to 81.8$ on MNIST,
$53.3 \to 23.1$ with large seed variance on CIFAR-10), consistent with
feedback alignment's known depth pathology. Claims about local wiring rules
in this paper are therefore scoped to shallow networks. Third, \emph{a
global scalar is not enough at any depth we tested}: \rmhebb{} stays far
below every vector-signal rule and degrades further with depth. Mask-budget
robustness at two hidden layers is shown in Figure~\ref{fig:sparsity}.

\paragraph{The observation that led to \fw{}.} The \textsc{readout} control
is the quiet protagonist: training \emph{only the top} of a frozen random
network already recovers $94.9\%$ on MNIST and $45.6\%$ on CIFAR-10---most
of the accuracy of full training, for a small fraction of its backward
cost. If cheap top-only updates carry that much of the signal, full-depth
credit assignment is needed only occasionally, and the question becomes
\emph{when}: the scheduling problem \fw{} answers.

\begin{figure}[t]
\centering
\includegraphics[width=\textwidth]{exp1_locality}
\caption{Learning curves for the wiring rules at two hidden layers (mean
over 3 seeds; bands $\pm 1$ std). Colored curves learn only the mask over
frozen random weights; gray lines mark the controls. At this depth the
vector-signal rules (\ep{}, \dfa{}) finish within two points of dense
training or above it; Table~\ref{tab:exp1} shows \dfa{}'s parity does not
survive depth 4.}
\label{fig:exp1curves}
\end{figure}

\begin{figure}[t]
\centering
\includegraphics[width=0.5\textwidth]{exp1_sparsity}
\caption{Mask-budget robustness on MNIST (2 hidden layers): both search
rules degrade gracefully down to $10\%$ of connections kept, where the local
rule still matches the trained-readout control.}
\label{fig:sparsity}
\end{figure}

\subsection{Two failure modes worth documenting}
\label{sec:exp1-stability}

\paragraph{The sign factor in straight-through mask gradients.} Selecting the
top-$k$ mask by $|s|$ while passing the straight-through gradient to $s$
directly drops the $\partial|s|/\partial s = \operatorname{sign}(s)$ factor;
every negative-score synapse then moves in the wrong direction, and learning
sits at chance. Applying $|\cdot|$ \emph{before} the straight-through
operator (so autodiff supplies the sign), or equivalently multiplying manual
updates by $\operatorname{sign}(s)$, restores learning entirely. We flag this
because the failure is silent: losses hover near initialization and no error
is raised.

\paragraph{Hebbian runaway and its stabilization.} A naive three-factor rule
diverges here: with ReLU activity the eligibility $a_i x_j$ is non-negative,
so reinforced co-activations grow activations, which grow eligibilities,
until logits saturate (magnitudes ${>}10^2$ within 200 steps), the softmax
collapses to one-hot, exploration dies, and accuracy pins at chance.
Batch-centered (covariance) eligibilities, RMS-normalized per-layer updates,
and a loss-based rather than sampled-action reward stabilize it; the
stabilized rule learns but remains far below the vector-signal rules.
